\documentclass[10pt,a4paper]{amsart}
\usepackage{amsmath,amssymb,amsthm,graphicx}
% Retained preamble for the Stacks Project chapter divisors.tex (Divisors).
% The chapter file pinned-original-source.tex carries only \input{preamble} on its first line, so no
% preamble text lives inside the chapter itself.  The Stacks Project uses one single shared preamble for
% every chapter; the snapshot stages it verbatim in sources/stacks-chow/preamble.tex.  The block below is
% copied mechanically, line for line, from that file's lines 14-18 (the Xy-pic setup) and lines 154-217
% (the native theorem declarations together with the mathematical macros).  Nothing was added, removed or
% redefined.  The wrapper class and the amsmath/amssymb/amsthm/graphicx base packages are supplied by the
% assembler, so the class/\externaldocument part of the shared preamble is not repeated here.
% For commutative diagrams we use Xy-pic
\usepackage[all]{xy}

% We use 2cell for 2-commutative diagrams.
\xyoption{2cell}
% Theorem environments.
%
\theoremstyle{plain}
\newtheorem{theorem}[subsection]{Theorem}
\newtheorem{proposition}[subsection]{Proposition}
\newtheorem{lemma}[subsection]{Lemma}

\theoremstyle{definition}
\newtheorem{definition}[subsection]{Definition}
\newtheorem{example}[subsection]{Example}
\newtheorem{exercise}[subsection]{Exercise}
\newtheorem{situation}[subsection]{Situation}

\theoremstyle{remark}
\newtheorem{remark}[subsection]{Remark}
\newtheorem{remarks}[subsection]{Remarks}

\numberwithin{equation}{subsection}

% Macros
%
\def\lim{\mathop{\mathrm{lim}}\nolimits}
\def\colim{\mathop{\mathrm{colim}}\nolimits}
\def\Spec{\mathop{\mathrm{Spec}}}
\def\Hom{\mathop{\mathrm{Hom}}\nolimits}
\def\Ext{\mathop{\mathrm{Ext}}\nolimits}
\def\SheafHom{\mathop{\mathcal{H}\!\mathit{om}}\nolimits}
\def\SheafExt{\mathop{\mathcal{E}\!\mathit{xt}}\nolimits}
\def\Sch{\mathit{Sch}}
\def\Mor{\mathop{\mathrm{Mor}}\nolimits}
\def\Ob{\mathop{\mathrm{Ob}}\nolimits}
\def\Sh{\mathop{\mathit{Sh}}\nolimits}
\def\NL{\mathop{N\!L}\nolimits}
\def\CH{\mathop{\mathrm{CH}}\nolimits}
\def\proetale{{pro\text{-}\acute{e}tale}}
\def\etale{{\acute{e}tale}}
\def\QCoh{\mathit{QCoh}}
\def\Ker{\mathop{\mathrm{Ker}}}
\def\Im{\mathop{\mathrm{Im}}}
\def\Coker{\mathop{\mathrm{Coker}}}
\def\Coim{\mathop{\mathrm{Coim}}}

% Boxtimes
%
\DeclareMathSymbol{\boxtimes}{\mathbin}{AMSa}{"02}

%
% Macros for moduli stacks/spaces
%
\def\QCohstack{\mathcal{QC}\!\mathit{oh}}
\def\Cohstack{\mathcal{C}\!\mathit{oh}}
\def\Spacesstack{\mathcal{S}\!\mathit{paces}}
\def\Quotfunctor{\mathrm{Quot}}
\def\Hilbfunctor{\mathrm{Hilb}}
\def\Curvesstack{\mathcal{C}\!\mathit{urves}}
\def\Polarizedstack{\mathcal{P}\!\mathit{olarized}}
\def\Complexesstack{\mathcal{C}\!\mathit{omplexes}}
% \Pic is the operator that assigns to X its picard group, usage \Pic(X)
% \Picardstack_{X/B} denotes the Picard stack of X over B
% \Picardfunctor_{X/B} denotes the Picard functor of X over B
\def\Pic{\mathop{\mathrm{Pic}}\nolimits}
\def\Picardstack{\mathcal{P}\!\mathit{ic}}
\def\Picardfunctor{\mathrm{Pic}}
\def\Deformationcategory{\mathcal{D}\!\mathit{ef}}

% --- appended 2026-10-12 (Chief): the project own definitions of the reference/slogan/history
% display environments, copied mechanically from the shared Stacks preamble lines 8-12
% (\usepackage{verbatim} plus the three \newenvironment lines, which the chapter preamble keeps
% verbatim from the project preamble). Nothing added, removed or redefined: these environments
% carry no mathematical content and the project itself typesets them as comments. Without these
% three lines a counted statement whose own text uses them fails with Environment slogan undefined.
\usepackage{verbatim}
\newenvironment{reference}{\comment}{\endcomment}
\newenvironment{slogan}{\comment}{\endcomment}
\newenvironment{history}{\comment}{\endcomment}
\usepackage{cleveref}
\crefname{originaldefinition}{original definition}{original definitions}
\Crefname{originaldefinition}{Original definition}{Original definitions}
\makeatletter
\newlabel{definition-blow-up}{{01OG}{source}{}{sourcepointer-1}{}}
\newlabel{definition-blow-up@cref}{{[originaldefinition][0][]01OG}{[1][0][]source}{}{}{}}
\makeatother
\crefname{originallemma}{original lemma}{original lemmas}
\Crefname{originallemma}{Original lemma}{Original lemmas}
\makeatletter
\newlabel{lemma-universal-property-blowing-up}{{0806}{source}{}{sourcepointer-2}{}}
\newlabel{lemma-universal-property-blowing-up@cref}{{[originallemma][0][]0806}{[1][0][]source}{}{}{}}
\makeatother
\title{A point of a blowing up lies in the affine chart of a local equation exactly when that equation cuts out the exceptional divisor there}
\author{The Stacks Project Authors}
\date{stacks-project:divisors.tex:local-snapshot-20261012; GNU FDL 1.2 or later}
\begin{document}
\maketitle
\noindent\textit{Editorial note.} All delivered mathematical text is the author's own, copied from the pinned Stacks Project chapter file pinned-original-source.tex; no statement or proof step was rewritten, reordered or removed.\par\medskip
\noindent\textit{Editorial note.} Following the multi-claim convention for this source, source\_url is the source-level URL https://stacks.math.columbia.edu/ and the source label lemma-characterize-affine-blowup is carried in claim\_key and here rather than in a per-claim tag URL.\par\medskip
\noindent\textit{Editorial note.} The preamble is the chapter's shared Stacks Project preamble, retained as sources/stacks-divisors/retained-preamble.tex: a line-for-line excerpt of the shared Stacks preamble staged with this snapshot. No macro was invented or redefined, and this card adds no preamble text of its own.\par\medskip
\noindent\textit{Editorial note.} Statement span 8469-8478 and proof span 8480-8503 come from the halfA pool card operations/pool-cards/stacks-divisors-halfA.json and were re-read in the pinned chapter file before this card was written.\par\medskip
\noindent\textit{Editorial note.} The proof uses the author's notation for the blowup algebra and writes a displayed map from the affine chart's coordinate ring into the local ring at the point; the display is preserved unchanged.\par\medskip
\noindent\textit{Editorial note.} 2 cross-references occur inside the counted spans; they are listed in external\_references as navigation pointers only.\par\medskip
\noindent\textit{Allowed standard prerequisites (editorial).}\begin{itemize}
\item The affine description of a blowing up by blowup algebras, with the exceptional divisor over the affine chart cut out by the image of the chosen element. Used at proof lines 8480-8483 and 8496-8498. Gives one direction immediately and supplies the description of the blowup algebra used to define the comparison map.
\item The universal property of blowing up, cited from the same chapter. Cited at proof lines 8501-8503. Upgrades the ring homomorphism built from the local divisibility into a morphism over the blowing up, which exhibits the point inside the affine chart.
\end{itemize}\medskip
\ifdefined\hypertarget\hypertarget{sourcepointer-1}{}\fi
\noindent\textit{Source reference.} Original definition 01OG: focused-seeds/AGP/sources/stacks-snapshot-20261011/tags.txt, source lines 2193--2193. Same-chapter reference used in the retained author text at source line(s) 8496. It is supplied only as a navigation pointer so that the retained text compiles; the counted item delivers the author's own statement and the author's complete proof, and no part of the referenced item is copied in or substituted for the counted argument.\par\medskip
\ifdefined\hypertarget\hypertarget{sourcepointer-2}{}\fi
\noindent\textit{Source reference.} Original lemma 0806: focused-seeds/AGP/sources/stacks-snapshot-20261011/tags.txt, source lines 10171--10171. Same-chapter reference used in the retained author text at source line(s) 8500, 8513, 8607, 8615, 8816. It is supplied only as a navigation pointer so that the retained text compiles; the counted item delivers the author's own statement and the author's complete proof, and no part of the referenced item is copied in or substituted for the counted argument.\par\medskip
\begin{lemma}
\label{lemma-characterize-affine-blowup}
Let $b : X' \to X$ be the blowing up of the scheme $X$ along a closed
subscheme $Z$. Let $U = \Spec(A)$ be an affine open of $X$ and let
$I \subset A$ be the ideal corresponding to $Z \cap U$.
Let $a \in I$ and let $x' \in X'$ be a point mapping to a point of $U$.
Then $x'$ is a point of the affine open $U' = \Spec(A[\frac{I}{a}])$
if and only if the image of $a$ in $\mathcal{O}_{X', x'}$ cuts
out the exceptional divisor.
\end{lemma}

\begin{proof}
Since the exceptional divisor over $U'$ is cut out by the image of
$a$ in $A' = A[\frac{I}{a}]$ one direction is clear. Conversely, assume
that the image of $a$ in $\mathcal{O}_{X', x'}$ cuts out $E$.
Since every element of $I$ maps to an element of the ideal
defining $E$ over $b^{-1}(U)$ we see that elements of $I$ become
divisible by $a$ in $\mathcal{O}_{X', x'}$. Thus for $f \in I^n$
we can write $f = \psi(f) a^n$ for some $\psi(f) \in \mathcal{O}_{X', x'}$.
Observe that since $a$ maps to a nonzerodivisor of $\mathcal{O}_{X', x'}$
the element $\psi(f)$ is uniquely characterized by this. Then we
define
$$
A' \longrightarrow \mathcal{O}_{X', x'},\quad
f/a^n \longmapsto \psi(f)
$$
Here we use the description of blowup algebras given following
Algebra, Definition \ref{definition-blow-up}. The uniqueness mentioned
above shows that this is an $A$-algebra homomorphism.
This gives a morphism $\Spec(\mathcal{O}_{X', x"}) \to \Spec(A') = U'$.
By the universal property of blowing up
(Lemma \ref{lemma-universal-property-blowing-up})
this is a morphism over
$X'$, which of course implies that $x' \in U'$.
\end{proof}

\end{document}
